\begin{table}[H]
\centering
\caption{Model Sensitivity: Demand Shock Employment and Welfare Across Calibrations}
\begin{threeparttable}
\small
\begin{tabular}{lcc}
\toprule
Scenario & $\Delta E_{48}$ (\%) & CE Welfare Loss (\%) \\
\midrule
Baseline (SMM) & $-0.9$ & 0.92 \\
Low scarring ($\lambda$=0.06) & $-0.9$ & 0.86 \\
High scarring ($\lambda$=0.18) & $-0.9$ & 0.89 \\
No scarring ($\lambda$=0) & $-0.9$ & 0.84 \\
No participation exit ($\chi$=0) & $-0.8$ & 0.93 \\
High participation exit ($\chi$=0.016) & $-1.1$ & 0.91 \\
Low matching efficiency (A=0.45) & $-0.7$ & 0.84 \\
High matching efficiency (A=0.75) & $-0.7$ & 0.92 \\
Fast mean reversion ($\rho$=0.98) & $-0.6$ & 0.56 \\
Slow mean reversion ($\rho$=0.995) & $-1.2$ & 1.41 \\
\bottomrule
\end{tabular}
\begin{tablenotes}[flushleft]
\small
\item \textit{Notes:} Each row reports the 48-month employment decline and CE welfare loss under an alternative calibration of one parameter. All other parameters are held at their SMM-estimated values. The demand shock follows an AR(1) process. Welfare is computed over a 600-month simulation horizon.
\end{tablenotes}
\end{threeparttable}
\label{tab:app_model_sensitivity}
\end{table}
